\medskip

\begin{enumerate}
	\item

On a $ \vect{\mathrm{AB}}  \begin{pmatrix}	-1 - (-2)\\ 3-0 \\ 0-2 \end{pmatrix} = \begin{pmatrix}	1\\ 3 \\ -2 \end{pmatrix}$
	et $ \vect{\mathrm{AC}}  \begin{pmatrix}	1 - (-2)\\ -1-0 \\ 2-2 \end{pmatrix} = \begin{pmatrix}	3\\ -1 \\ 0 \end{pmatrix}$

$\dfrac{3}{1}\neq \dfrac{-1}{3}$ donc les vecteurs  $ \vect{\mathrm{AB}} $ et  $ \vect{\mathrm{AC}} $ ne sont pas colinéaires donc les points A, B et C ne sont pas alignés
		\item
		\begin{enumerate}
			\item
$\vect{n} \cdot \vect{\mathrm{AB}}= 1 \times 1 + 3 \times 3 + 5 \times -2 = 1 +9 -10=0$

$\vect{n} \cdot \vect{\mathrm{AC}}= 1 \times 3 + 3 \times -1 + 5 \times 0 =3-3+0=0$

Le vecteur $\vect{n}\begin{pmatrix}1 \\ 3 \\ 5\end{pmatrix}$ est donc orthogonal à deux vecteurs non colinéaires du plan (ABC) donc $\vect{n}$ est orthogonal au plan (ABC).

		\item  Le vecteur $\vect{n}$ est orthogonal au plan (ABC), c'est donc un vecteur normal du plan (ABC).

Une équation cartésienne du plan (ABC) est donc de la forme :

$x+3y+5z+d=0$ avec $d\in\R$

De plus, le point A appartient au plan (ABC) donc ses coordonnées verifient l'équation du plan. On a donc :

$-2 +3\times 0 + 5\times 2 +d=0 \Longleftrightarrow -2+10+d=0 \Longleftrightarrow d=-8$

Une équation cartésienne du plan (ABC) est donc	$x+3 y+5 z-8=0$.

		\item  $x_D+3 y_D+5 z_D-8=0+3\times 0 +5\times 3 -8=15-8=7\neq 0 $

Donc le point D n'appartient pas au plan (ABC) d'où  les points A, B, C et D ne sont pas coplanaires.
		\end{enumerate}

	\item

		\begin{enumerate}
			\item  Une équation paramétrique de $\mathcal{D}_{1}$ est :  $\begin{cases} x=t \\ y=3 t \\ z=3+5 t  \end{cases} $ avec $t \in \R$

D'une part : pour $t=0$, on a $\begin{cases} x=0 \\ y=3 \times0 =0  \\ z=3+5 \times 0 =3  \end{cases} $.

			On reconnait les coordonnées du point D, donc D appartient à la droite $\mathcal{D}_{1}$.

D'autre part : un vecteur directeur de $\mathcal{D}_{1}$ est donc le vecteur $\begin{pmatrix}	1\\ 3 \\ 5 \end{pmatrix}$ c'est à dire le vecteur $\vect{n}$ qui est orthogonal au plan(ABC).

Donc la droite $\mathcal{D}_{1}$ est orthogonal au plan (ABC).


Donc  la droite $\mathcal{D}_{1}$ est bien la hauteur du tétraèdre ABCD issue de D.

		\item  Pour déterminer les points éventuels d'intersection des droites $\mathcal{D}_{1}$ et $\mathcal{D}_{2}$  résolvons le système (S) :  $\begin{cases} t=1+3s \\ 3 t=-1-5s \\ 3+5 t=2-6s  \end{cases} $

(S) $\Longleftrightarrow \begin{cases} t =1+3s \\ 3  \times (1+3s)=-1-5s \\ 3+5 \times ( 1+3s)=2-6s  \end{cases}  \Longleftrightarrow  \begin{cases} t =1+3s \\ 3+9s=-1-5s \\ 3+5+15s=2-6s  \end{cases}$

(S) $\Longleftrightarrow  \begin{cases} t =1+3s \\ 14s=-4\\ 21s=-6  \end{cases} \Longleftrightarrow  \begin{cases} t =1+3s = 1 + 3\times \dfrac{-2}{7}= \dfrac{1}{7} \\ s=-\dfrac{4}{14}=-\dfrac{2}{7}\\ s=-\dfrac{6}{21} =-\dfrac{2}{7} \end{cases} $

Le système admet une unique solution donc les droites $\mathcal{D}_{1}$ et $\mathcal{D}_{2}$   sont sécantes.


Remplaçons $t$ par $\dfrac{1}{7}$ dans l'équation paramétrique de $\mathcal{D}_{1}$

			  $\begin{cases} x= \dfrac{1}{7}\\ y=3 \times \dfrac{1}{7} = \dfrac{3}{7} \\ z=3+5 \times \dfrac{1}{7}  = \dfrac{26}{7} \end{cases} $

Les droites $\mathcal{D}_{1}$ et $\mathcal{D}_{2}$ sont donc sécantes et  les coordonnées de leur point d'intersection sont $\left( \dfrac{1}{7}~;~\dfrac{3}{7}~;~\dfrac{26}{7}\right)$.
\end{enumerate}
	\item
		\begin{enumerate}
		\item  Soit H le projeté orthogonal  du point D sur le plan (ABC).

H est donc l'intersection du plan (ABC) et de la hauteur issue de D dans le tétraèdre ABCD c'est à dire la droite $\mathcal{D}_{1}$.

Les coordonnées de H vérifient donc l'équation cartésienne du plan (ABC) et l'équation paramétrique de la droite  $\mathcal{D}_{1}$.

$t+3 \times (3t) +5 \times (3+5t) -8=0 \Longleftrightarrow t+9t+15+25t-8=0 \Longleftrightarrow 35t = -7 \Longleftrightarrow t= -\dfrac{7}{35} = -\dfrac{1}{5}$

Remplaçons $t$ par $-\dfrac{1}{5}$ dans l'équation paramétrique de $\mathcal{D}_{1}$

$\begin{cases} x=- \dfrac{1}{5}\\ y=3 \times \dfrac{-1}{5} = -\dfrac{3}{5} \\ z=3+5 \times \dfrac{-1}{5}  = \dfrac{10}{5}=2 \end{cases} $

Le projeté orthogonal  du point D sur le plan (ABC) est le point H de coordonnées H$\left( - \dfrac{1}{5}~;~- \dfrac{3}{5}~;~2\right)$.

	\item La distance du point D au plan (ABC) est égale à la longueur DH car H est le projeté orthogonal  du point D sur le plan (ABC).

$\mathrm{DH}^2 = {  { \left( - \dfrac{1}{5} - 0  \right) } ^2  + { \left( - \dfrac{3}{5} - 0  \right) } ^2  + (2-3)^2 }= {  \dfrac{1}{25}  +  \dfrac{9}{25}  + 1} ={ \dfrac{35}{25}}$

DH$=\sqrt{ \dfrac{35}{25}}= \sqrt{ \dfrac{7}{5}} = \sqrt{\dfrac{14}{10}} = \sqrt{1,4}$ ou $\dfrac{\sqrt{35}}{5} \approx 1,183$ soit 1,18 au centième près.
	\end{enumerate}
\end{enumerate}
